% !PRINTED 385
\weberpage{385}

\weq{5}{\omega_n' + 1 = \frac{1}{Q_n}\left\{ Q_n - Q_{n-1} - \frac{(-1)^n}{Q_n\left(\omega' - \dfrac{P_n}{Q_n}\right)} \right\}.}

Wenn nun $n$ hinlänglich gross angenommen wird, so unterscheiden sich
\[
\omega' - \frac{P_n}{Q_n}, \quad \omega' - \frac{P_{n-1}}{Q_{n-1}}
\]
% [sic] The source prints $\omega'-\omega=2y\sqrt d$; from (1)--(2) of §122
% the left side is actually $-2y\sqrt d$.  The proof here uses only non-vanishing.
beliebig wenig von der nicht verschwindenden Differenz $\omega' - \omega = 2y\sqrt{d}$, wonach (4) zeigt, dass $\omega_n'$ negativ ist. Aus (5) aber folgt, dass $\omega_n' + 1$ positiv ist, da $Q_n - Q_{n-1}$ als positive ganze Zahl mindestens gleich 1 ist, und
\[
\frac{1}{Q_n\left(\omega' - \dfrac{P_n}{Q_n}\right)}
\]
für ein hinlänglich grosses $n$ beliebig klein wird; und dadurch ist der Beweis des Satzes 1. geführt.

Da die conjugirte Zahl von der conjugirten $\omega'$ die ursprüngliche Zahl $\omega$ ist, so folgt aus der Definition der reducirten Zahlen, \gesperrt{dass $\omega$ und $-1:\omega'$ gleichzeitig reducirt sind.}

\gesperrt{2. Ist $\omega$ reducirt und für ein ganzzahliges $\alpha$}

\weq{6}{\omega = \alpha + \frac{1}{\omega_1},}

\gesperrt{so ist $\omega_1$ dann und nur dann reducirt, wenn $\alpha$ die grösste in $\omega$ enthaltene ganze Zahl ist, wenn also}

\weq{7}{\alpha < \omega < \alpha + 1.}

Ist die Bedingung (7) erfüllt, so ist $\omega_1$ positiv und grösser als 1, und

\weq{8}{\omega_1' = \frac{1}{\omega' - \alpha}}

ist, da $\omega'$ negativ ist, ein negativer echter Bruch; also ist $\omega_1$ reducirt. Dagegen kann $\omega_1$ nicht reducirt sein, wenn $\alpha > \omega$ ist, da dann $\omega_1$ negativ wäre, oder wenn $\alpha + 1 < \omega$ wäre, da sonst $\omega_1$ kleiner als 1 wäre.

Setzt man (8) in die Form

\weq{9}{-\frac{1}{\omega_1'} = \alpha + \frac{1}{-\dfrac{1}{\omega'}},}
%CONT>>
